\section*{Chapitre \ref{ch:b3:complete} --- Espaces complets~:
Baire, Ascoli, Stone--Weierstrass}

\begin{solution}{exo:b3:complete:1}
(a) Pour $m \geq n$, $f_n - f_m$ s'annule hors d'un intervalle
de longueur $\frac1n$ et est bornée par $1$~: $\norm{f_n -
f_m}_1 \leq \frac1n \to 0$~: de Cauchy. Si $f_n \to f$ en
$\norm\cdot_1$ avec $f$ \omterm{def:b3:topology:continuity}{continue}~: pour $\alpha < \frac12$
fixé, $\int_0^{\alpha}\abs f = \lim \int_0^\alpha\abs{f -
f_n + f_n} \leq \lim\bigl(\norm{f - f_n}_1 + 0\bigr) = 0$
($f_n \equiv 0$ là pour $n$ grand), donc $f \equiv 0$ sur
$[0, \frac12)$ (continuité)~; de même $f \equiv 1$ sur
$(\frac12, 1]$~: aucune fonction \omterm{def:b3:topology:continuity}{continue} ne fait les deux.
Donc l'espace est incomplet --- le \omterm{thm:b3:complete:completion}{complété} est $L^1$, construit
dans le~\cref{ch:b3:lp}.

(b) Soit $(x_n)$ de Cauchy~; choisir $n_1 < n_2 < \cdots$ avec
$\norm{x_{n_{k+1}} - x_{n_k}} \leq 2^{-k}$. La série
$\sum_k (x_{n_{k+1}} - x_{n_k})$ converge absolument, donc
converge~; ses sommes partielles sont $x_{n_{k+1}} - x_{n_1}$,
donc $(x_{n_k})$ converge, et une suite de Cauchy admettant
une sous-suite convergente converge.
\end{solution}

\begin{solution}{exo:b3:complete:2}
(a) Supposer $E$ \omterm{def:b3:complete:complete}{complet} avec \omterm{def:b3:topology:basis}{base} \omterm{def:b3:galois:algebraic}{algébrique} $(e_n)_{n\in\N}$
et soit $F_n = \operatorname{Vect}(e_1, \dots, e_n)$~: fermé
(les sous-espaces de dimension finie sont \omterm{def:b3:complete:complete}{complets}, donc
fermés --- deuxième année), d'intérieur vide~: si $B(x, r)
\subseteq F_n$, prendre $v \notin F_n$~; alors $x +
\frac{r}{2\norm v}v \in B(x,r) \setminus F_n$, absurde. Or $E
= \bigcup F_n$ (tout vecteur est combinaison finie)~:
contredit Baire (\cref{thm:b3:complete:baire}). L'espace
$\R[X]$ a la \omterm{def:b3:topology:basis}{base} dénombrable $(X^n)$, donc aucune norme ne le
rend \omterm{def:b3:complete:complete}{complet}.

(b) Fixer $\delta > 0$ et une boule ouverte non vide $B$~; on
trouve dans $B$ un point de $\Omega_\delta = \{x :
\operatorname{osc}_xf < \delta\}$, où $\operatorname{osc}_xf
= \inf_{V \ni x}\operatorname{diam} f(V)$. ($\Omega_\delta$
est \omterm{def:b3:topology:topology}{ouvert}~: si $\operatorname{diam}f(V) < \delta$ pour un
\omterm{def:b3:topology:topology}{ouvert} $V \ni x$, tout $y \in V$ a une oscillation $<
\delta$.) Les ensembles
\[
F_N = \bigl\{x : \abs{f_n(x) - f_m(x)} \leq \tfrac\delta3\
\ \forall\, n, m \geq N \bigr\}
\]
sont fermés (intersections d'images réciproques de fermés) et
recouvrent $\R$ (la convergence ponctuelle rend $(f_n(x))$ de
Cauchy). Appliquer Baire dans le \omterm{def:b3:complete:complete}{complet} $\bar B$~: quelque
$F_N \cap \bar B$ contient une boule $B' = B(x_0, \rho)$. En
faisant $m \to \infty$~: $\abs{f_N - f} \leq \frac\delta3$ sur
$B'$. Par continuité de $f_N$ en $x_0$, rétrécir en $B'' \ni
x_0$ où $\abs{f_N - f_N(x_0)} \leq \frac\delta3$~; alors pour
$x \in B''$,
\[
\abs{f(x) - f(x_0)} \leq \abs{f - f_N}(x) + \abs{f_N(x) -
f_N(x_0)} + \abs{f_N - f}(x_0) \leq \delta:
\]
$\operatorname{diam} f(B'') \leq 2\delta$, donc $x_0 \in
\Omega_{3\delta} \cap B$. Ainsi chaque $\Omega_\delta$ est
\omterm{def:b3:topology:topology}{ouvert} et \omterm{def:b3:topology:interior}{dense}~; $\bigcap_k\Omega_{1/k}$ --- l'ensemble des
points de continuité --- est \omterm{def:b3:topology:interior}{dense} par Baire.
\end{solution}

\begin{solution}{exo:b3:complete:3}
(a) $f(x) = \frac12(x + \frac ax)$ envoie $[\sqrt a, \infty)$
dans lui-même (AM--GM~: $f(x) \geq \sqrt{x \cdot \frac ax} =
\sqrt a$), et $f'(x) = \frac12(1 - \frac a{x^2}) \in [0,
\frac12)$ là~: une $\frac12$-contraction lipschitzienne d'un
ensemble fermé (\omterm{def:b3:complete:complete}{complet}). Point fixe~: $x = f(x) \iff x^2 =
a$~: $x^* = \sqrt a$. Taux (\cref{thm:b3:complete:banach})~:
$d(x_n, \sqrt a) \leq 2^{-n+1}\,d(x_1, x_0)$. Pour $a = 2$,
$x_0 = 2$~: $d(x_1, x_0) = \frac12$, donc $n = 40$ garantit
$2^{-40} < 10^{-12}$. (En réalité la méthode de Newton
converge quadratiquement~: une poignée d'itérations
suffisent~; l'estimation de contraction est pessimiste mais
gratuite.)

(b) $f_m(x) = m + e\sin x$ est $e$-lipschitzienne avec $e < 1$
sur le \omterm{def:b3:complete:complete}{complet} $\R$~: unique point fixe $x(m)$. Pour deux
paramètres~:
\[
\abs{x(m) - x(m')} = \abs{f_m(x(m)) - f_{m'}(x(m'))}
\leq \abs{m - m'} + e\abs{x(m) - x(m')},
\]
donc $\abs{x(m) - x(m')} \leq \frac{\abs{m - m'}}{1 - e}$~:
même lipschitzienne en $m$.
\end{solution}

\begin{solution}{exo:b3:complete:4}
(a) Soit $x^*$ l'unique point fixe de $f^p$. Alors
$f^p(f(x^*)) = f(f^p(x^*)) = f(x^*)$~: $f(x^*)$ est un point
fixe de $f^p$, donc $f(x^*) = x^*$. Un point fixe de $f$ en
est un de $f^p$~: l'unicité se transfère. Pour $T$~: par
induction $\abs{T^pf(x)} \leq \norm f_\infty x^p/p!$ (chaque
intégration ajoute un facteur $\frac xk$), donc $\norm{T^p}
\leq a^p/p!$. L'application $S(u) = g + \lambda Tu$ satisfait
$S^p(u) - S^p(v) = \lambda^pT^p(u - v)$, de norme $\leq
\abs\lambda^pa^p/p!\,\norm{u - v} \to 0$~: quelque $S^p$ est
une contraction, et $S$ a un unique point fixe~: l'équation de
Volterra $u = g + \lambda Tu$ est uniquement solvable pour
\emph{tout} $\lambda$.

(b) $\varphi(x) = d(x, f(x))$ est \omterm{def:b3:topology:continuity}{continue} sur le compact
$K$~: elle atteint son minimum en quelque $x_0$. Si $f(x_0)
\neq x_0$, alors $\varphi(f(x_0)) = d(f(x_0), f^2(x_0)) <
d(x_0, f(x_0)) = \min\varphi$~: absurde. Unicité~: deux points
fixes $x \ne y$ donneraient $d(x,y) = d(f(x), f(y)) <
d(x,y)$. Sans compacité~: $f(x) = x + \frac1x$ sur $[1,
\infty)$ satisfait, pour $x < y$, $f(y) - f(x) = (y -
x)\bigl(1 - \frac1{xy}\bigr) < y - x$, pourtant $f(x) > x$
toujours~: pas de point fixe --- la décroissance stricte des
distances est plus faible qu'un facteur de contraction
uniforme.
\end{solution}

\begin{solution}{exo:b3:complete:5}
(a) L'ensemble $A = \{g = h\}$ est l'image réciproque de la
diagonale $\Delta_Y$ sous $x \mapsto (g(x), h(x))$, \omterm{def:b3:topology:continuity}{continue}~;
$\Delta_Y$ est fermée car $Y$ est \omterm{def:b3:topology:hausdorff}{séparé} (pour $y_1 \neq
y_2$, des \omterm{def:b3:topology:topology}{voisinages} \omterm{def:b3:topology:topology}{ouverts} disjoints donnent un pavé \omterm{def:b3:topology:topology}{ouvert}
autour de $(y_1, y_2)$ disjoint de la diagonale, donc le
complémentaire de $\Delta_Y$ est \omterm{def:b3:topology:topology}{ouvert})~: donc $A$ est
fermé, contient un \omterm{def:b3:topology:interior}{dense}, égale $X$. Utilisé~: unicité dans
\cref{thm:b3:complete:extension}, donc dans l'unicité des
\omterm{thm:b3:complete:completion}{complétés}.

(b) $f$, une isométrie, est uniformément \omterm{def:b3:topology:continuity}{continue}~: elle
s'étend en $F\colon X \to Y$
(\cref{thm:b3:complete:extension}), encore isométrique (la
relation $d(F(x), F(x')) = d(x, x')$ tient sur un \omterm{def:b3:topology:interior}{dense} de
paires et les deux côtés sont continus). De même
$f^{-1}\colon f(D) \to X$ s'étend en $G \colon Y \to X$. Le
composé $G \circ F$ est continu et fixe le \omterm{def:b3:topology:interior}{dense} $D$~: c'est
$\mathrm{id}_X$ (partie (a))~; symétriquement $F \circ G =
\mathrm{id}_Y$. Donc $F$ est une bijection isométrique.
Unicité des \omterm{thm:b3:complete:completion}{complétés}~: appliquer ceci à $D = X$ densément
assis dans deux \omterm{thm:b3:complete:completion}{complétés}.
\end{solution}
\begin{solution}{exo:b3:complete:6}
$\{\sin(nx)\}$~: bornée ponctuellement par $1$~; \emph{pas}
\omterm{def:b3:complete:equicontinuous}{équicontinue}~: en $x = 0$, $\sin(n\cdot\frac{\pi}{2n}) = 1$
avec $\frac\pi{2n} \to 0$, violant tout $\delta$ commun pour
$\varepsilon = \frac12$. Pas relativement compacte (nécessité
d'Ascoli, \cref{thm:b3:complete:ascoli}).

$\{x^n\}$~: bornée~; pas \omterm{def:b3:complete:equicontinuous}{équicontinue} en $1$~: $1 - (1 -
\delta)^n \to 1$ quand $n \to \infty$ pour $\delta$ fixé. Pas
relativement compacte --- cohérent, sa limite ponctuelle est
discontinue, donc aucune sous-suite ne converge uniformément.

$\{\norm f_\infty \leq 1, \norm{f'}_\infty \leq 1\}$~:
l'inégalité des accroissements finis rend la famille
$1$-lipschitzienne, donc \omterm{def:b3:complete:equicontinuous}{équicontinue}~; bornée~: relativement
compacte par Ascoli. (Pas compacte~: elle n'est pas fermée ---
les limites uniformes n'ont pas besoin d'être $\mathcal C^1$~;
son \omterm{def:b3:topology:interior}{adhérence} est les fonctions $1$-lipschitziennes de norme
$\leq 1$.)

$\{\operatorname{Lip}(f) \leq 1, f(0) = 0\}$~: \omterm{def:b3:complete:equicontinuous}{équicontinue}~;
bornée ponctuellement ($\abs{f(x)} \leq x \leq 1$)~; et fermée
sous les limites uniformes (l'inégalité de Lipschitz et la
valeur en $0$ passent aux limites)~: compacte.
\end{solution}

\begin{solution}{exo:b3:complete:7}
(a) Pour $\norm f_\infty \leq 1$~: $\abs{Tf(x)} \leq
\norm k_\infty$, et
\[
\abs{Tf(x) - Tf(x')} \leq \int_0^1\abs{k(x,y) -
k(x',y)}\,\dd y \leq \omega_k\bigl(\abs{x - x'}\bigr),
\]
où $\omega_k$ est un \omterm{def:b3:modules:module}{module} de continuité uniforme de $k$ sur
le carré compact (Heine)~: l'image de la boule unité est
uniformément bornée et \omterm{def:b3:complete:equicontinuous}{équicontinue}, donc relativement
compacte (Ascoli). Par linéarité tout ensemble borné a une
image relativement compacte~: $T$ est un opérateur compact.

(b) L'énoncé de sous-suite est la définition de compacité
relative appliquée à $\{Tf_n\}$. Si $T$ était bijective
d'inverse \omterm{def:b3:topology:continuity}{continue}, alors $T(B(0,1)) \supseteq B(0,
\varepsilon)$ pour quelque $\varepsilon > 0$ ($T^{-1}$
\omterm{def:b3:topology:continuity}{continue} en $0$)~; la boule fermée $\bar B(0,\varepsilon)$,
partie fermée du compact $\overline{T(B(0,1))}$, serait
compacte --- impossible dans le $\mathcal C(\intcc01)$ de
dimension infinie par le théorème de Riesz (deuxième année).
\end{solution}

\begin{solution}{exo:b3:complete:8}
(a) Dini~: voir \cref{lem:b3:complete:dini} --- l'argument de
recouvrement. (b) Faux~: la bosse mobile $f_n(x) = \max(0, 1
- \abs{nx - 1})$ tend vers $0$ ponctuellement (pour $x > 0$,
$f_n(x) = 0$ dès que $n > 2/x$~; $f_n(0) = 0$) mais
$\norm{f_n}_\infty = 1$. (c) Faux~: $f_n(x) = x^n$ décroît
vers la discontinue $\mathbf 1_{\{1\}}$~;
$\sup_{[0,1]}\abs{f_n - f} = \sup_{[0,1)} x^n = 1$. (d)
Faux~: $x^n \downarrow 0$ ponctuellement sur le non-compact
$\intoo01$, avec $\sup_{(0,1)}x^n = 1$. Chaque hypothèse de
Dini est nécessaire.
\end{solution}

\begin{solution}{exo:b3:complete:9}
(a) Par linéarité $\int_0^1 fP = 0$ pour tout polynôme $P$.
Choisir $P_n \to f$ uniformément
(\cref{cor:b3:complete:weierstrass})~: $\int_0^1 f^2 =
\lim\int_0^1 fP_n = 0$, et la \omterm{def:b3:topology:continuity}{continue} $f^2 \geq 0$
d'intégrale nulle s'annule identiquement.

(b) Sur $\intcc01$~: les polynômes en $x^2$ forment une
algèbre avec constantes, séparant les points ($x \mapsto x^2$
est injective sur $\intcc01$)~: \omterm{def:b3:topology:interior}{dense} par Stone--Weierstrass.
Sur $\intcc{-1}1$~: $x^2$ prend des valeurs égales en $\pm
x$, et de même tout polynôme en $x^2$~: une limite uniforme
d'une telle suite est une fonction paire. Si des fonctions
paires $g_n$ convergeaient uniformément vers l'identité, alors
$x = \lim g_n(x) = \lim g_n(-x) = -x$ pour tout $x$~:
absurde --- pas \omterm{def:b3:topology:interior}{dense}. L'hypothèse de séparation échoue sur
les paires $\{x, -x\}$.

(c) Non. Pour $f$ dans l'algèbre $\mathcal A$ engendrée par
les constantes et $\eu^{\iu t}$ --- combinaisons linéaires de
$\eu^{\iu nt}$, $n \geq 0$ --- on a $\Lambda(f) =
\int_0^{2\pi}f(t)\,\eu^{\iu t}\,\dd t = 0$ (chaque
$\int_0^{2\pi}\eu^{\iu(n+1)t}\dd t = 0$). $\Lambda$ est
\omterm{def:b3:topology:continuity}{continue} pour $\norm\cdot_\infty$ ($\abs{\Lambda(f)}\leq
2\pi\norm f_\infty$), donc $\Lambda$ s'annule sur
$\bar{\mathcal A}$~; mais $\Lambda(\eu^{-\iu t}) = 2\pi \neq
0$~: $\eu^{-\iu t} \notin \bar{\mathcal A}$.
(Stone--Weierstrass ne s'applique pas~: $\mathcal A$ n'est
pas stable par conjugaison --- et l'obstruction est
précisément celle que la théorie des fonctions holomorphes
systématisera dans le~\cref{ch:b3:holomorphic}.)
\end{solution}

\begin{solution}{exo:b3:complete:10}
$F_n = \bigcap_{f \in \mathcal F}\{x : \abs{f(x)} \leq n\}$
est une intersection de fermés~: fermé. La bornitude
ponctuelle donne $X = \bigcup_n F_n$. Baire
(\cref{thm:b3:complete:baire}) fournit $n_0$ avec $U =
\mathring F_{n_0} \neq \varnothing$~: sur $U$, $\abs f \leq
n_0$ pour toute $f \in \mathcal F$ simultanément.
\end{solution}

\begin{solution}{exo:b3:complete:11}
(a) Une suite de Cauchy pour $\norm\cdot_{\mathcal C^1}$ est
uniformément de Cauchy avec ses dérivées~: $f_n \to f$ et
$f_n' \to g$ uniformément, avec $f, g$ \omterm{def:b3:topology:continuity}{continues}. En passant
à la limite (la convergence uniforme le permet sous
l'intégrale) dans $f_n(x) = f_n(0) + \int_0^xf_n'(t)\,\dd t$
on obtient $f(x) = f(0) + \int_0^xg$~: $f$ est $\mathcal C^1$
avec $f' = g$, et $\norm{f_n - f}_{\mathcal C^1} \to 0$.
\omterm{def:b3:complete:complete}{Complet}.

(b) $h(x) = \abs{x - \frac12}$ est une limite uniforme de
fonctions $\mathcal C^1$, p.\,ex.\ $h_n(x) = \sqrt{(x -
\frac12)^2 + \frac1n}$ ($\abs{h_n - h} \leq \frac1{\sqrt n}$
par la borne de quantité conjuguée), pourtant $h \notin E$~:
la norme sup sur $E$ n'est pas complète --- son \omterm{thm:b3:complete:completion}{complété} est
$\mathcal C(\intcc01)$.

(c) $f_n(x) = \sin(2\pi nx)$ a $\norm{f_n}_\infty = 1$ et
$\norm{f_n'}_\infty = 2\pi n \to \infty$~: aucun $C$
n'existe. (Conceptuellement~: si la dérivation était bornée
pour la norme sup, les deux normes de (a)--(b) seraient
équivalentes, rendant $(E, \norm\cdot_\infty)$ \omterm{def:b3:complete:complete}{complet} ---
contredisant (b). Cette non-bornitude est exactement pourquoi
les fonctions \omterm{def:b3:topology:continuity}{continues} génériques peuvent n'être dérivables
nulle part, \cref{thm:b3:complete:nowherediff}.)
\end{solution}

\begin{solution}{exo:b3:complete:12}
Fixer $\varepsilon > 0$. Chaque $F_N$ est une intersection
sur $n \geq N$ d'images réciproques du fermé
$\intcc{-\varepsilon}\varepsilon$ sous la \omterm{def:b3:topology:continuity}{continue} $x \mapsto
f(nx)$~: fermé. L'hypothèse dit que tout $x > 0$ est dans
quelque $F_N$. Par Baire appliqué dans le \omterm{def:b3:complete:complete}{complet} $\intcc ab$
(n'importe quels $0 < a < b$), quelque $F_N$ est \omterm{def:b3:topology:interior}{dense} dans
un sous-intervalle~; étant fermé, il contient un intervalle
$\intcc{a'}{b'}$ avec $0 < a' < b'$. Alors pour tout $n \geq
\max(N, \frac{a'}{b' - a'})$, les intervalles
$\intcc{na'}{nb'}$ et $\intcc{(n+1)a'}{(n+1)b'}$ se
chevauchent ($nb' \geq (n+1)a'$), donc
\[
\bigcup_{n \geq n_0}\intcc{na'}{nb'} \supseteq
\intco{n_0a'}{+\infty},
\]
et tout $t \geq n_0a'$ s'écrit $t = nx$ avec $n \geq N$, $x
\in \intcc{a'}{b'} \subseteq F_N$~: $\abs{f(t)} \leq
\varepsilon$. Donc $\limsup_{t\to\infty}\abs f \leq
\varepsilon$ pour tout $\varepsilon$~: $f \to 0$.
L'hypothèse complète est nécessaire car $F_N$ doit
\emph{recouvrir} une valeur d'intervalle de $x$ --- avec
seulement les $x$ rationnels la réunion des $F_N$ est
dénombrable et Baire ne donne rien~; il existe d'ailleurs des
contre-exemples continus s'annulant le long de tous les
rayons rationnels mais pas à l'infini.
\end{solution}
\begin{solution}{pb:b3:complete:1}
\textbf{1.} Induction sur $k$~: si $\norm{\varphi_n(t_k) -
x_0} \leq M(t_k - t_0) \leq MT \leq b$, le point $(t_k,
\varphi_n(t_k))$ est dans $R$, donc la pente $f(t_k,
\varphi_n(t_k))$ est définie, de norme $\leq M$~; alors pour
$t \in [t_k, t_{k+1}]$, $\norm{\varphi_n(t) - x_0} \leq
\norm{\varphi_n(t_k) - x_0} + M(t - t_k) \leq M(t - t_0)
\leq MT \leq b$.

\textbf{2.} Chaque morceau affine a une pente de norme $\leq
M$~; une fonction affine par morceaux à pentes bornées par
$M$ est $M$-lipschitzienne (chaîner via les nœuds).

\textbf{3.} La famille $(\varphi_n)$ est bornée
ponctuellement (valeurs dans $\bar B(x_0, b)$) et
\omterm{def:b3:complete:equicontinuous}{équicontinue} (constante de Lipschitz commune $M$)~: Ascoli
(\cref{thm:b3:complete:ascoli}) extrait $\varphi_{n_j} \to
\varphi$ uniformément sur $[t_0, t_0 + T]$. Les bornes
passent à la limite~: $\varphi$ est $M$-lipschitzienne,
$\varphi(t_0) = x_0$.

\textbf{4.} $R$ est compact et $f$ \omterm{def:b3:topology:continuity}{continue}~: uniformément
\omterm{def:b3:topology:continuity}{continue} (Heine, \cref{cor:b3:topology:heineborel})~; soit
$\delta(\varepsilon)$ un \omterm{def:b3:modules:module}{module}. Pour $t$ non nœud $\in
(t_k, t_{k+1})$~: $\varphi_n'(t) = f(t_k, \varphi_n(t_k))$,
et les deux points d'évaluation de $f$ diffèrent de $\abs{t -
t_k} \leq T/n$ en temps et $\norm{\varphi_n(t) -
\varphi_n(t_k)} \leq MT/n$ en espace. Pour $n \geq n_0$ avec
$(1 + M)T/n_0 < \delta(\varepsilon)$~: $\norm{\Delta_n(t)}
\leq \varepsilon$.

\textbf{5.} Sur chaque $[t_k, t_{k+1}]$, $\varphi_n$ est
affine, donc $\varphi_n(t_{k+1}) - \varphi_n(t_k) =
\int_{t_k}^{t_{k+1}} \varphi_n'(s)\dd s$, la dérivée étant
la pente constante~; en sommant sur les morceaux (et en
coupant le dernier en $t$)~: $\varphi_n(t) = x_0 +
\int_{t_0}^t\varphi_n'(s)\,\dd s$. En écrivant $\varphi_n' =
f(s, \varphi_n(s)) + \Delta_n(s)$ (intégrandes \omterm{def:b3:topology:continuity}{continues} par
morceaux, un nombre fini de sauts) on obtient l'affichage.

\textbf{6.} Étant donné $\varepsilon$~: pour $j$ grand,
$\norm{\varphi_{n_j} - \varphi}_\infty <
\delta(\varepsilon)$, donc $\norm{f(s, \varphi_{n_j}(s)) -
f(s, \varphi(s))} \leq \varepsilon$ pour tout $s$~:
convergence uniforme des intégrandes, et $\int_{t_0}^t
f(s,\varphi_{n_j}(s))\dd s \to \int_{t_0}^tf(s,
\varphi(s))\dd s$ uniformément en $t$. Aussi
$\norm{\int_{t_0}^t\Delta_{n_j}} \leq
T\sup\norm{\Delta_{n_j}} \to 0$ (question~4). En passant à
la limite dans l'identité de la question~5~: $\varphi(t) =
x_0 + \int_{t_0}^tf(s, \varphi(s))\,\dd s$.

\textbf{7.} L'intégrande $s \mapsto f(s, \varphi(s))$ est
\omterm{def:b3:topology:continuity}{continue}, donc le second membre est $\mathcal C^1$ en $t$ de
dérivée $f(t, \varphi(t))$~: $\varphi$ résout le problème de
Cauchy sur $[t_0, t_0 + T]$. Pour la moitié gauche, poser
$g(t, x) = -f(2t_0 - t, x)$, \omterm{def:b3:topology:continuity}{continue} sur le rectangle
réfléchi avec la même borne $M$~; une solution $\psi$ de $y'
= g(t,y)$, $y(t_0) = x_0$ sur $[t_0, t_0 + T]$ donne
$\varphi(t) = \psi(2t_0 - t)$ résolvant l'équation originale
sur $[t_0 - T, t_0]$~; les deux moitiés se recollent en une
solution $\mathcal C^1$ (les deux dérivées unilatérales en
$t_0$ valent $f(t_0, x_0)$). --- \emph{Théorème de Peano}~:
une $f$ \omterm{def:b3:topology:continuity}{continue} admet une solution locale par toute
condition initiale.

\textbf{8.} La continuité est claire. Lipschitz près de $0$
échoue~: $\abs{f(h) - f(0)} = 2\sqrt h$, et $2\sqrt h \leq L
h$ est faux pour $h < 4/L^2$.

\textbf{9.} Pour $t \leq c$~: $x_c \equiv 0$ résout. Pour $t
\geq c$~: $x_c'(t) = 2(t - c) = 2\sqrt{(t-c)^2} =
2\sqrt{\abs{x_c}}$. En $t = c$ les dérivées unilatérales
valent toutes deux $0$~: $x_c$ est $\mathcal C^1$ et résout
globalement, avec $x_c(0) = 0$ pour tout $c \geq 0$ --- avec
$x \equiv 0$, un continuum de solutions par l'origine.

\textbf{10.} L'itération de Picard pose $\Phi(x)(t) = x_0 +
\int_0^t f(x(s))\dd s$ et a besoin de $\norm{\Phi(x) -
\Phi(y)} \leq k\norm{x - y}$ avec $k < 1$ sur une boule
convenable --- ce qui suit d'une borne de Lipschitz sur $f$,
transférée sous l'intégrale. Ici $\sqrt\cdot$ n'admet aucune
borne de Lipschitz près de $0$, et aucun choix d'intervalle
ou de boule ne le répare. Et en effet aucune preuve
d'unicité ne pourrait réussir~: l'unicité est \emph{fausse}
(question~9).

\textbf{11.} Dans $c_0$, les vecteurs unitaires $e_n$
satisfont $\norm{e_n - e_m}_\infty = 1$ pour $n \neq m$~:
aucune sous-suite n'est de Cauchy, donc la boule unité
fermée n'est pas compacte. L'étape qui casse est l'extraction
(question~3)~: Ascoli pour $\mathcal C([t_0, t_0+T], E)$
exige que les valeurs vivent dans un espace où les ensembles
bornés sont relativement compacts --- vrai dans $\R^d$
(Bolzano--Weierstrass), faux dans $c_0$.

\textbf{12.} \emph{Peano}~: $f$ \omterm{def:b3:topology:continuity}{continue} au \omterm{def:b3:topology:topology}{voisinage} de
$(t_0, x_0)$ dans $\R\times\R^d$ $\Rightarrow$ il existe une
solution $\mathcal C^1$ de $x' = f(t,x)$, $x(t_0) = x_0$, sur
quelque $[t_0 - T, t_0 + T]$. \emph{Cauchy--Lipschitz}
(\cref{ch:b3:ode})~: si de plus $f$ est localement
lipschitzienne en la variable $x$, la solution est
\emph{unique}. La paire $(x' = 2\sqrt{\abs x},\ x(0) = 0)$
sépare les deux théorèmes.

\textbf{13.} $\omega(r) = Lr$~: $\int_0^1\frac{\dd r}{Lr} =
+\infty$~: qualifie. $\omega(r) = r\log\frac1r$~:
$\int\frac{\dd r}{r\log(1/r)} = \bigl[-\log\log\frac1r\bigr]
\to +\infty$~: qualifie --- pourtant $\omega(r)/r =
\log\frac1r \to \infty$~: pas lipschitzien. $\omega(r) =
2\sqrt r$~: $\int_0^1\frac{\dd r}{2\sqrt r} = 1 < \infty$~:
échoue.

\textbf{14.} Soustraire les deux formes intégrales et prendre
les normes avec le \omterm{def:b3:modules:module}{module} d'Osgood~: $\delta(t) \leq
\delta(s) + \int_s^t\omega(\delta(v))\,\dd v$.

\textbf{15.} $\tau$ est bien défini ($\delta(t_0) = 0$) avec
$\delta(\tau) = 0$ et $\delta > 0$ sur $\intoc\tau{t_1}$.
Fixer $s \in \intoo\tau{t_1}$. Alors $u$ est $\mathcal C^1$,
$u(s) = \delta(s) > 0$, $u \geq \delta$ sur $[s, t_1]$, et
$u' = \omega(\delta) \leq \omega(u)$. Diviser et intégrer~:
\[
\int_{\delta(s)}^{u(t_1)}\frac{\dd r}{\omega(r)}
\leq t_1 - s \leq t_1 - \tau .
\]
Le second membre est borné, tandis que le premier $\to
+\infty$ quand $s \downarrow \tau$ (divergence en $0$)~:
contradiction. Donc $\delta \equiv 0$~: unicité.

\textbf{16.} Lipschitz est $\omega = Lr$~: unicité
retrouvée. Pour $x' = x\log\frac1{\abs x}$~: le second membre
satisfait le \omterm{def:b3:modules:module}{module} d'Osgood $\omega(r) = r\log\frac1r$ près
de $0$~: solutions uniques. Pour $\omega(r) = 2\sqrt r$~:
$\int_0^{h^2}\frac{\dd r}{2\sqrt r} = h$ --- le «~budget
d'Osgood~» pour grimper de $0$ à la hauteur $h^2$ est
exactement le temps $h$, et en effet $x_c(c + h) = h^2$.

\textbf{17.} $G(t) = \int_{t_0}^t(Le(s) + \eta(s))\dd s$
vérifie $G' \leq LG + \sup\eta$. Alors
$\bigl(\eu^{-Lt}(G + \tfrac{\sup\eta}L)\bigr)' \leq 0$~: $e(t)
\leq G(t) \leq \frac{\sup\eta}L(\eu^{L(t-t_0)} - 1)$.

\textbf{18.} Défaut quantitatif~: $\norm{\Delta_n} \leq L'T/n
+ LMT/n$. Soustraire les identités intégrales et utiliser la
borne de Lipschitz avec la question~17~: la borne $O(1/n)$
énoncée. Même sans taux~: toute sous-suite de
$(\varphi_n)$ a une sous-sous-suite convergeant (Ascoli +
partie~II) vers \emph{une} solution, que l'unicité force à
être $\varphi$~: une suite dont toutes les sous-suites ont
des sous-sous-suites de même limite converge.

\textbf{19.} Depuis $\varphi_n(t_k) = 0$~: la pente $0$ donne
$\varphi_n \equiv 0$. Depuis $\varepsilon > 0$~: sur
$[\varepsilon', \infty)$ la fonction $2\sqrt x$ est
lipschitzienne, donc Euler converge vers $(t +
\sqrt\varepsilon)^2$~; faire $\varepsilon \to 0$ atterrit sur
$t^2$, pas sur $0$.

\textbf{20.} Non vide~: partie~II. Toute solution satisfait
$\norm{x'} \leq M$~: $\mathcal S$ est uniformément
$M$-lipschitzienne et bornée. Fermée~: passer à la limite
dans la forme intégrale. Ascoli~: $\mathcal S$ est compacte.

\textbf{21.} Toute solution est croissante avec $x(0) = 0$.
Soit $c = \sup\{t : x(t) = 0\}$~; pour $t > c$, $(\sqrt x)' =
1$, donc $x = x_c$. La famille se réduit à $\{x_c : c \in
[0,1]\}$ (avec $x_1 = 0$), image \omterm{def:b3:topology:continuity}{continue} du compact \omterm{def:b3:topology:connected}{connexe}
$[0,1]$~: compacte et \omterm{def:b3:topology:connected}{connexe}.

\textbf{22.} L'évaluation $\operatorname{ev}_t$ est \omterm{def:b3:topology:continuity}{continue},
donc $\mathcal S(t)$ est compacte et \omterm{def:b3:topology:connected}{connexe}. Pour
l'exemple~: $x_c(t) = (t - c)_+^2$ balaye $\intcc0{t^2}$.

\textbf{23.} Soustraire les formes intégrales~: $e(t) \leq d
+ \int_{t_0}^t L\,e$~; Grönwall donne $e(t) \leq
d\,\eu^{L(t-t_0)}$. Optimalité sur $x' = Lx$. Unicité pour $d
= 0$. Le flot est $\eu^{LT}$-lipschitzien en la condition
initiale.

\textbf{24.} $\Omega$ est une bijection croissante~; son
inverse $g$ résout $g' = \omega(g)$ avec $g(0^+) = 0$.
Étendue par $0$ pour $t \leq 0$, $g$ est une solution
$\mathcal C^1$ distincte de $0$~: quand l'intégrale
converge, l'unicité échoue. Pour $\omega(r) = 2\sqrt r$~:
$g(t) = t^2$, et les translations donnent la famille $x_c$.

\textbf{25.} Avec pas $h = \frac1n$~: $\varphi_n(1) = (1 +
\frac1n)^n$. Le développement $n\log(1+\frac1n) = 1 -
\frac1{2n} + O(n^{-2})$ donne $(1+\frac1n)^n =
\eu(1-\frac1{2n}+O(n^{-2}))$~: erreur $\frac{\eu}{2n} +
O(n^{-2})$. Numériquement $n = 10$~: $1.1^{10} \approx
2.59374$ contre $\eu \approx 2.71828$, erreur $0.12454$ vs
$\frac{\eu}{20} \approx 0.13591$.


En détail pour Osgood et les taux (questions~13--18)~: le
\omterm{def:b3:modules:module}{module} $\omega(r) = Lr$ donne l'intégrale divergente
$\int_0\frac{\dd r}{Lr} = +\infty$, tandis que $2\sqrt r$
donne $\int_0^1\frac{\dd r}{2\sqrt r} = 1 < \infty$, ce qui
explique exactement pourquoi $x_c$ peut quitter $0$ en temps
fini $h$ pour atteindre la hauteur $h^2$. La preuve d'Osgood
construit une majorante $u \geq \delta$ d'équation
différentielle $u' \leq \omega(u)$~; diviser par $\omega(u)$
et intégrer produit une borne inférieure de l'intégrale de
$\frac{\dd r}{\omega(r)}$ par $t_1 - \tau$, qui explose
quand le bord inférieur $\delta(s)$ tend vers $0$. Grönwall
intégral suit en dérivant $\eu^{-Lt}G$~: le facteur
exponentiel est optimal, atteint par $x' = Lx$. Avec des
données lipschitziennes, le défaut d'Euler est
$O(\frac1n)$, et Grönwall le convertit en erreur $O(\frac1n)$
sur la solution~; l'unicité promeut en outre la convergence
sous-séquentielle (Ascoli) en convergence de toute la suite.

Pour l'entonnoir de solutions (questions~19--22)~: le schéma
d'Euler, démarré exactement en $0$, reste en $0$ (pente
nulle)~; démarré en $\varepsilon > 0$, il converge vers $(t +
\sqrt\varepsilon)^2$, qui tend vers $t^2$ quand
$\varepsilon \to 0$ --- sensibilité numérique de la
non-unicité. L'ensemble $\mathcal S$ de toutes les solutions
est fermé (passage à la limite dans la forme intégrale),
borné et équicontinu (borne $M$), donc compact par Ascoli.
Pour $x' = 2\sqrt{\abs x}$, $\mathcal S = \{x_c : c \in
[0,1]\}$ est l'image \omterm{def:b3:topology:continuity}{continue} d'un segment, donc compacte et
\omterm{def:b3:topology:connected}{connexe}~; la section $\mathcal S(t) = \intcc0{t^2}$ montre
que tout état intermédiaire est réalisé.

Pour l'optimalité et le calcul explicite (questions~23--25)~:
la dépendance lipschitzienne en la condition initiale suit
de Grönwall avec terme constant $\norm{x_0 - y_0}$. Quand
$\int_0\frac{\dd r}{\omega(r)}$ converge, l'inverse de
$\Omega(x) = \int_0^x\frac{\dd r}{\omega(r)}$ fournit une
solution non nulle partant de $0$~: la divergence est la
frontière exacte. Enfin $(1+\frac1n)^n =
\eu(1-\frac1{2n}+O(n^{-2}))$ confirme le taux $O(\frac1n)$
avec constante exacte $\frac\eu2$, vérifiée numériquement
pour $n = 10$.

\end{solution}
